\documentclass[a4paper]{arrowhead}

\usepackage[yyyymmdd]{datetime}
\usepackage{etoolbox}
\usepackage[utf8]{inputenc}
\usepackage[T1]{fontenc}
\usepackage{multirow}
\usepackage{longtable}
\usepackage{tikz}
\usepackage{adjustbox}
\AtBeginEnvironment{longtable}{\small}
\BeforeBeginEnvironment{tikzpicture}{\begin{adjustbox}{max width=\textwidth}}
\AfterEndEnvironment{tikzpicture}{\end{adjustbox}}
\usetikzlibrary{positioning,shapes.multipart,arrows.meta}

\renewcommand{\dateseparator}{-}
\setlength{\emergencystretch}{3em}
\lstset{language=HTTP,breakatwhitespace=false,columns=fullflexible}

%% Special references
\newcommand{\mref}[2][]{{\textcolor{ArrowheadPurple}{\hyperref[sec:model:#2]{\ifstrempty{#1}{#2}{#1}}}}}
\newcommand{\pref}[1]{{\textcolor{ArrowheadGrey}{\hyperref[sec:model:primitives:#1]{#1}}}}
%%

\begin{document}
\sloppy

%% Arrowhead Document Properties
\ArrowheadTitle{SimpleStoreServiceOrchestration}
\ArrowheadType{System Description}
\ArrowheadTypeShort{SysD}
\ArrowheadVersion{5.2.2}
\ArrowheadDate{2026-11-07}
\ArrowheadAuthor{Jerker Delsing}
\ArrowheadStatus{RELEASE CANDIDATE}
\ArrowheadContact{jerker.delsing@ltu.se}
\ArrowheadFooter{\href{www.arrowhead.eu}{www.arrowhead.eu}}
\ArrowheadSetup
%%

%% Front Page
\begin{center}
  \vspace*{1cm}
  \huge{\arrowtitle}

  \vspace*{0.2cm}
  \LARGE{\arrowtype}
  \vspace*{1cm}

  \vspace*{\fill}

  % Front Page Image

  \vspace*{1cm}
  \vspace*{\fill}

  % Front Page Abstract
  \begin{abstract}
    This document describes the SimpleStoreServiceOrchestration core system of the Eclipse Arrowhead 5.2.2 framework, as implemented in Java Spring in \texttt{ah5-core-java-spring/serviceorchestration-simple}, including the local extensions of this workspace.
  \end{abstract}

  \vspace*{1cm}
\end{center}

\newpage
%%

%% Table of Contents
\tableofcontents
\newpage
%%

\section{Overview}
\label{sec:overview}

This document describes the SimpleStoreServiceOrchestration system. Simple\allowbreak{}Store\allowbreak{}Service\allowbreak{}Orchestration orchestrates from a static store: an operator records which service instance each consumer shall use, with a priority, and the orchestrator returns these rules on request. It offers the same \texttt{service\allowbreak{}Orchestration} service as Dynamic\allowbreak{}Service\allowbreak{}Orchestration, plus management of the store.

\begin{figure}[ht!]
  \centering
  \begin{tikzpicture}[font=\small, >=Stealth]
    \node[draw=ArrowheadBlue, very thick, minimum width=4.2cm, minimum height=2.2cm, align=center, fill=ArrowheadBlue!10] (s)
      {\guillemotleft{}System\guillemotright{}\\\textbf{SimpleStoreServiceOrchestration}};
    \node[draw=ArrowheadTeal, rectangle, align=left, left=1.3cm of s, text width=4.6cm, inner sep=5pt] (p)
      {\textbf{Produced services}\\ \texttt{service\allowbreak{}Orchestration}\\ \texttt{service\allowbreak{}Orchestration\allowbreak{}Simple\allowbreak{}Store\allowbreak{}Management}\\ \texttt{service\allowbreak{}Orchestration\allowbreak{}Push\allowbreak{}Management}\\ \texttt{service\allowbreak{}Orchestration\allowbreak{}History\allowbreak{}Management}\\ \texttt{generalManagement}\\ \texttt{sysMonitor}};
    \node[draw=ArrowheadOrange, rectangle, align=left, right=1.3cm of s, text width=4.6cm, inner sep=5pt] (c)
      {\textbf{Consumed}\\ AMQP broker (SysMonitor)\\ Authentication\\ Blacklist\\ ConsumerAuthorization\\ ServiceRegistry\\ Subscribers};
    \draw[->, thick, ArrowheadTeal] (s) -- (p);
    \draw[->, thick, ArrowheadOrange] (s) -- (c);
  \end{tikzpicture}
  \caption{SysML-style block view of SimpleStoreServiceOrchestration with its produced and consumed services.}
  \label{fig:system}
\end{figure}

The rest of this document is organized as follows.
In Section \ref{sec:prior_art}, we reference major prior art capabilities of the system.
In Section \ref{sec:use}, we describe the intended usage of the system.
In Section \ref{sec:properties}, we describe fundamental properties provided by the system.
In Section \ref{sec:delimitations}, we describe delimitations of capabilities of the system.
In Section \ref{sec:services}, we describe the abstract service functions consumed or produced by the system.
In Section \ref{sec:security}, we describe the security capabilities of the system.

\newpage

\subsection{Significant Prior Art}
\label{sec:prior_art}

Successor of the store-based orchestration of the Arrowhead 4.x Orchestrator. In Arrowhead 5 it is a separate, deliberately simple system.

\subsection{How This System Is Meant to Be Used}
\label{sec:use}

An operator creates store entries with \texttt{service\allowbreak{}Orchestration\allowbreak{}Simple\allowbreak{}Store\allowbreak{}Management.\allowbreak{}create}. A consumer pulls with the service definition (and optional preferred providers and flags \texttt{MATCHMAKING}, \texttt{ONLY\_\allowbreak{}PREFERRED}) and receives the instance identifiers in priority order. It then looks up the interfaces of the chosen instance in the ServiceRegistry. Under \texttt{service.\allowbreak{}discovery.\allowbreak{}policy=\allowbreak{}restricted}, the default, that lookup only works when the consumer is listed in \texttt{service.\allowbreak{}discovery.\allowbreak{}direct.\allowbreak{}access}; otherwise the instance is silently missing from the result. In the test clouds it runs alongside Dynamic\allowbreak{}Service\allowbreak{}Orchestration; Java consumers pin the orchestrator by provider name.

\subsection{System functionalities and properties}
\label{sec:properties}

\subsubsection{Functional properties of the system}

\begin{itemize}
  \item Pull orchestration from the store in priority order, with preferred providers and matchmaking.
  \item Push orchestration with subscriptions and triggers, as Dynamic\allowbreak{}Service\allowbreak{}Orchestration.
  \item Management of store entries: create, query, change priorities, remove.
  \item Orchestration history.
\end{itemize}

\subsubsection{Configuration of system properties}

The system is configured with Spring properties (\texttt{application.properties}, overridden by environment variables in the container deployment). The properties that shape its behaviour are:

\noindent\begin{longtable}{| p{5.4cm} | p{3.4cm} | p{6.4cm} |} \hline
\rowcolor{gray!33} Property & Default & Meaning \\ \hline \endhead
\texttt{server.\allowbreak{}port} & 8442 & HTTPS port (port specification). The upstream image still carries 8465; an image built from this tree does not. \\ \hline
\texttt{authentication.\allowbreak{}policy} & declared & Test cloud: \texttt{outsourced}. \\ \hline
\texttt{arrowhead.\allowbreak{}security.\allowbreak{}profile} & secure & \texttt{development} selects the isolated development exception; never set in a deployment. \\ \hline
\texttt{management.\allowbreak{}policy} / \texttt{management.\allowbreak{}whitelist} & authorization & Test cloud: whitelist LocalCloudManager (PDE now uses DSO). \\ \hline
\texttt{enable.\allowbreak{}blacklist.\allowbreak{}filter} & false & Blacklist filter of requesters (test cloud: true). \\ \hline
\texttt{orchestration.\allowbreak{}history.\allowbreak{}max.\allowbreak{}age} & 15 & Days the history is kept. \\ \hline
\texttt{cleaner.\allowbreak{}job.\allowbreak{}interval} & 30000 & Cleaner interval in ms. \\ \hline
\texttt{push.\allowbreak{}orchestration.\allowbreak{}max.\allowbreak{}thread} & 5 & Push worker threads. \\ \hline
\texttt{normalization.\allowbreak{}mode} & extended & Name normalization. \\ \hline
\texttt{mqtt.\allowbreak{}api.\allowbreak{}enabled} & false & Offer the MQTT interfaces. \\ \hline
\texttt{sysmonitor.\allowbreak{}enabled} & false & Publish native sysMonitor observations (test cloud: true). \\ \hline
\texttt{arrowhead.\allowbreak{}ssl.\allowbreak{}crl-\allowbreak{}file} & (unset) & CRL checked for every client certificate (local). \\ \hline
\end{longtable}

Secrets (database, keystore and identity passwords, secret keys) have no usable defaults; in the test clouds they are per-cloud secrets mounted as files.

\subsubsection{Data stored by the system}

Orchestration store entries (consumer, service definition, instance, priority, creator), push subscriptions, orchestration jobs and log entries (database \texttt{ah\_\allowbreak{}serviceorchestration\_\allowbreak{}simple}).

\subsubsection{Non functional properties}

A pull is a single database query; no other system is called during orchestration, so it works while Consumer\allowbreak{}Authorization is unavailable.

\subsubsection{Stateful or stateless}

Stateful: store entries, subscriptions and history are persistent.

\subsection{Important Delimitations}
\label{sec:delimitations}

The store is not checked against the ServiceRegistry or Consumer\allowbreak{}Authorization; a rule can point to an expired instance or an unauthorized provider. Results contain no interfaces, metadata or tokens. Only the flags \texttt{MATCHMAKING} and \texttt{ONLY\_\allowbreak{}PREFERRED} are supported.

\newpage

\section{Services}
\label{sec:services}

This section describes the produced and consumed services. Each produced service is specified by an SD document and one IDD document per interface.

\subsection{Produced services}

\noindent\begin{longtable}{| p{5.3cm} | p{1.3cm} | p{2.6cm} | p{1.6cm} | p{3.6cm} |} \hline
\rowcolor{gray!33} Service & Version & Interfaces & SD & IDD \\ \hline \endhead
\texttt{service\allowbreak{}Orchestration} & 1.0.0 & HTTP(S), MQTT(S) & \cite{ssso-serviceOrchestration-sd} & \cite{ssso-serviceOrchestration-http-idd}, \cite{ssso-serviceOrchestration-mqtt-idd} \\ \hline
\texttt{service\allowbreak{}Orchestration\allowbreak{}Simple\allowbreak{}Store\allowbreak{}Management} & 1.0.0 & HTTP(S), MQTT(S) & \cite{ssso-serviceOrchestrationSimpleStoreManagement-sd} & \cite{ssso-serviceOrchestrationSimpleStoreManagement-http-idd}, \cite{ssso-serviceOrchestrationSimpleStoreManagement-mqtt-idd} \\ \hline
\texttt{service\allowbreak{}Orchestration\allowbreak{}Push\allowbreak{}Management} & 1.0.0 & HTTP(S), MQTT(S) & \cite{ssso-serviceOrchestrationPushManagement-sd} & \cite{ssso-serviceOrchestrationPushManagement-http-idd}, \cite{ssso-serviceOrchestrationPushManagement-mqtt-idd} \\ \hline
\texttt{service\allowbreak{}Orchestration\allowbreak{}History\allowbreak{}Management} & 1.0.0 & HTTP(S), MQTT(S) & \cite{ssso-serviceOrchestrationHistoryManagement-sd} & \cite{ssso-serviceOrchestrationHistoryManagement-http-idd}, \cite{ssso-serviceOrchestrationHistoryManagement-mqtt-idd} \\ \hline
\texttt{generalManagement} & 1.0.0 & HTTP(S), MQTT(S) & \cite{ssso-generalManagement-sd} & \cite{ssso-generalManagement-http-idd}, \cite{ssso-generalManagement-mqtt-idd} \\ \hline
\texttt{sysMonitor} & 1.0.0 & AMQP(S) & \cite{ssso-sysMonitor-sd} & \cite{ssso-sysMonitor-amqp-idd} \\ \hline
\end{longtable}

The HTTP interfaces are always offered, the MQTT interfaces with \texttt{mqtt.api.enabled=true} and the AMQP \texttt{sysMonitor} interface with \texttt{sysmonitor.enabled=true}. Upstream 5.2.2 also answers \texttt{GET <base>/monitor/echo} on some systems; this endpoint is not registered as a service.

\subsection{Consumed services}

\noindent\begin{longtable}{| >{\raggedright\arraybackslash}p{2.9cm} | >{\raggedright\arraybackslash}p{3.1cm} | >{\raggedright\arraybackslash}p{2.5cm} | p{3.6cm} | >{\raggedright\arraybackslash}p{2.9cm} |} \hline
\rowcolor{gray!33} Provider & Service & Operations & Use & Condition \\ \hline \endhead
ServiceRegistry & serviceDiscovery & register, revoke, lookup & Registration; lookup of core services. & always \\ \hline
Authentication & identity & identity-login, identity-verify & Own token; verification of requester tokens. & \texttt{authentication.\allowbreak{}policy=\allowbreak{}outsourced} \\ \hline
Consumer\allowbreak{}Authorization & authorization & verify & Management filter decisions. & \texttt{management.\allowbreak{}policy=\allowbreak{}authorization} \\ \hline
Blacklist & blacklist\allowbreak{}Discovery & check & Blacklist filter. & \texttt{enable.\allowbreak{}blacklist.\allowbreak{}filter=\allowbreak{}true} \\ \hline
Subscribers & notify interface & HTTP(S) request or MQTT(S) publish & Push orchestration results. & subscriptions \\ \hline
AMQP broker (SysMonitor) & sysMonitor (publisher) & echo, power, status & Own observations. & \texttt{sysmonitor.\allowbreak{}enabled=\allowbreak{}true} \\ \hline
\end{longtable}

Except for the bootstrap services (\texttt{http.collector.bootstrap.services}), which are looked up in the ServiceRegistry, the system reaches other systems' services through DynamicServiceOrchestration (\texttt{http.collector.mode=ORCHESTRATION}, results cached for \texttt{http.collector.orchestration.ttl} = 300 s); every such call needs a ConsumerAuthorization grant.

\newpage

\section{Security}
\label{sec:security}

The default is the Arrowhead secure mode with HTTPS and X.509 certificates following the Arrowhead certificate profile (\texttt{<System>.<Cloud>.<Organization>.arrowhead.eu}). Certificates not issued under the cloud CA are refused by the TLS handshake. The Arrowhead-test-cloud-java deployment runs the secure mode.

The packaged \texttt{application.properties} is a development sample that disables TLS and declares identities. The system refuses to start with it, or with any other effective configuration that sets \texttt{server.ssl.enabled=false}, \texttt{authentication.policy=declared}, \texttt{service.authorization.enabled=false} or, while monitoring is on, \texttt{sysmonitor.ssl.enabled=false}. The check runs before the system accepts a connection. Only an explicit \texttt{arrowhead.security.profile=development} lets such a system start, for an isolated test, and it then logs a non-conformance warning.

\subsection{Security model}

\begin{itemize}
  \item \textbf{Protocols:} HTTPS (TLS 1.3 and 1.2) with mandatory client certificates (\texttt{server.ssl.client-auth=need}); MQTT over TLS; AMQP over TLS for \texttt{sysMonitor}.
  \item \textbf{Data protection:} transport encryption by TLS; the database connection uses TLS with server verification in the test cloud. Secrets are kept out of the container environment (Compose file secrets).
  \item \textbf{System authentication:} by the authentication policy: \texttt{certificate} (client certificate CN), \texttt{outsourced} (identity token verified with Authentication), \texttt{internal} (Authentication only) or \texttt{declared} (insecure). In this workspace an identity token is accepted only together with the client certificate of the system it was issued to (token-to-certificate binding), and every client certificate is checked against the cloud CA's CRL when \texttt{arrowhead.ssl.crl-file} is set (fail closed).
  \item \textbf{Authorization of produced services:} management services are protected by the management filter (\texttt{management.policy}: \texttt{SYSOP\_ONLY}, \texttt{WHITELIST} or \texttt{AUTHORIZATION} via ConsumerAuthorization); an optional blacklist filter refuses blacklisted requesters. Sysop identities bypass the management filter.
  \item Because the store bypasses Consumer\allowbreak{}Authorization, operators must keep store entries consistent with the authorization policies; providers should still check the consumer's identity. Three combinations with the ServiceRegistry discovery policy are intended: (1) Dynamic\allowbreak{}Service\allowbreak{}Orchestration for ordinary consumers, the recommended default, where no consumer needs discovery access; (2) \texttt{restricted} with each SimpleStore consumer named in \texttt{service.\allowbreak{}discovery.\allowbreak{}direct.\allowbreak{}access}, which also lets those consumers discover every other service, so a provider must still check them; (3) \texttt{open}, only for an isolated cloud where every system may discover every service. A consumer that holds the access details in its local configuration needs no lookup and no discovery access.
\end{itemize}

For the Arrowhead certificate profile see github.com/eclipse-arrowhead/documentation.


\bibliographystyle{IEEEtran}
\bibliography{bibliography}

\newpage

\section{Revision History}
\subsection{Amendments}

\noindent\begin{tabularx}{\textwidth}{| p{1cm} | p{3cm} | p{2cm} | X | p{4cm} |} \hline
\rowcolor{gray!33} No. & Date & Version & Subject of Amendments & Author \\ \hline

1 & 2026-11-07 & \arrowversion & Document created from the Arrowhead LaTeX templates for the 5.2.2 Java Spring implementation, including the local extensions of this workspace. & Jerker Delsing \\ \hline
\end{tabularx}

\subsection{Quality Assurance}

\noindent\begin{tabularx}{\textwidth}{| p{1cm} | p{3cm} | p{2cm} | X |} \hline
\rowcolor{gray!33} No. & Date & Version & Approved by \\ \hline

1 & & \arrowversion  &  \\ \hline

\end{tabularx}

\end{document}
